% TOP_PROBLEM: {"problemId": "problem.dynamical-mordell-lang-conjecture", "canonicalTitle": "Dynamical Mordell–Lang Conjecture", "releaseRank": 251, "primaryDomain": "number_theory_arithmetic_geometry", "primaryDomainLabel": "Number theory & arithmetic geometry", "displayStatus": "open_with_solved_subcases", "releaseStatus": "open_with_solved_subcases"}
% !TeX program = lualatex
% ATTEMPT_STATUS: unresolved
% Research continuation by GPT_6_Astra_Ultra, 27 September 2026.
\documentclass[11pt]{article}
\usepackage[margin=25mm]{geometry}
\usepackage{fontspec}
\setmainfont{Latin Modern Roman}
\usepackage{amsmath,amssymb,mathtools}
\usepackage{unicode-math}
\setmathfont{Latin Modern Math}
\usepackage{xurl,hyperref}
\hypersetup{colorlinks=true,urlcolor=blue}
\setcounter{secnumdepth}{0}
\setlength{\parindent}{0pt}
\setlength{\parskip}{5pt}
\setlength{\emergencystretch}{3em}
\begin{document}
\section*{\texorpdfstring{Dynamical Mordell–Lang Conjecture}{Dynamical Mordell–Lang Conjecture}}\label{251}
\subsection{Definitions and mathematical statement}
For a morphism $f:X\to X$ of a quasiprojective variety over a characteristic-zero field $K$, a closed subvariety $V$, and $x\in X(K)$, define the return-time set
\[
 N_f(x,V)=\{n\in{\mathbb{Z}}_{\ge0}:f^n(x)\in V(K)\}.
\]
The conjecture says this is a finite union of sets $a+b{\mathbb{Z}}_{\ge0}$ with $a,b\ge0$. The case $b=0$ is a singleton; an empty union is allowed. Thus exceptional isolated returns and infinite periodic strings of returns are both permitted. Since $f$ is everywhere defined, all iterates exist. No dominance, smoothness, invertibility, or \"etaleness is assumed in the universal statement.
\par\smallskip\textit{Formulation and concept references:} \href{https://arxiv.org/pdf/2602.08608v2}{[S1]}.

\subsection{Short English statement}
Must the times at which an algebraic orbit visits a subvariety consist of finitely many arithmetic progressions and isolated times?
\subsection{Research attempt: exact interpolation and the ramification gap}
\textit{GPT-6 Astra Ultra, 27 September 2026. Status: UNRESOLVED.}

The useful target is an exact description of the set of times, not an upper bound on its density. I first prove a class of cases where each residue class of times has polynomial interpolation, then isolate what an analytic replacement would have to supply. The characteristic-zero assumption is tested by a counterexample in positive characteristic. These arguments do not remove the general obstruction caused by ramification.

\paragraph{A polynomial-orbit lemma.}
Suppose $f$ is an affine map on $\mathbb A^d_K$. Write its action in homogeneous affine coordinates as
\[
 \binom{f(z)}1=A\binom z1,
 \qquad A\in M_{d+1}(K).
\]
Assume that for some $q\ge1$ one has $A^q=I+N$ with $N^r=0$. For $0\le a<q$ and every integer $m\ge0$,
\[
 A^{a+qm}=A^a\sum_{j=0}^{r-1}\binom mjN^j.                 \tag{1}
\]
Characteristic zero makes each $\binom mj$ a polynomial in $m$ with coefficients in $K$. Thus the coordinates of $f^{a+qm}(x)$ are polynomial functions of $m$, of degree at most $r-1$.

Let $P_1,\ldots,P_s$ generate the ideal of $V$ in the affine coordinate ring. Substitution in (1) gives finitely many polynomials $Q_{i,a}(m)$. If all are identically zero, the entire progression $a+q\mathbb Z_{\ge0}$ is a return set. Otherwise one nonzero $Q_{i,a}$ has only finitely many roots, and the simultaneous zero set is finite. Taking the union over $a$ proves the conjectured conclusion for this class of affine maps and arbitrary closed $V$. The proof does not suppose that $V$ is linear or irreducible.

There is also an explicit bound on the exceptional part. If some nonzero $Q_{i,a}$ comes from an equation of degree $D_i$, then the number of exceptional times in that residue class is at most $D_i(r-1)$. This is a bound on distinct roots; multiple roots do not create additional times. A constant nonzero equation gives no returns. The bound is only meaningful in the residue classes where an equation remains nonzero: one cannot apply it to a progression lying identically in $V$.

For example, let $f(u,v)=(u+1,-v)$ and $x=(0,1)$. For
\[
 V:\quad (u-3)(v-1)=0,
\]
we have $f^n(x)=(n,(-1)^n)$, so the return set is exactly
\[
 2\mathbb Z_{\ge0}\ \cup\ \{3\}.
\]
This checks both features allowed in the formulation: an infinite progression and an isolated time that does not lie on it. It also shows why replacing the statement by ``eventually either always or never'' would be too restrictive unless residue classes were first separated.

\paragraph{The finite-orbit case does not require interpolation.}
If $f^a(x)=f^b(x)$ with $a<b$, determinism gives $f^{n+b-a}(x)=f^n(x)$ for every $n\ge a$. Membership in $V$ is therefore periodic after time $a$. The finitely many earlier returns are singletons, and each returning position in a period supplies a progression. This elementary argument applies to all morphisms in the statement, including constant maps and maps with a preperiodic starting point. It does not require injectivity: repetition in a forward orbit is enough.

The difficult case consequently has an infinite orbit. In particular, enumerating a long finite segment without seeing a repetition proves neither eventual periodicity of the orbit nor a structural theorem for its intersection with $V$. The orbit can be infinite even when every return set in a particular family of targets is finite.

\paragraph{A precise analytic sufficient condition.}
Here is the step a $p$-adic method would need. Suppose that after embedding the relevant finitely generated field of definition into a suitable complete characteristic-zero nonarchimedean field, there exist $q\ge1$, $a_0\ge0$, and analytic maps
\[
 F_a:\mathbb Z_p\longrightarrow U_a\subseteq X,
 \qquad F_a(m)=f^{a_0+a+qm}(x),\quad 0\le a<q,
                                                               \tag{2}
\]
where each $U_a$ admits finitely many analytic equations $P_{i,a}=0$ cutting out $V$ along the image. Assume specifically that the compositions $P_{i,a}\circ F_a$ are restricted convergent power series on the closed unit disc. Then the conjecture follows for this orbit.

Indeed, a nonzero such series has finitely many zeros in $\mathbb Z_p$, by the nonarchimedean zero theorem used in the $p$-adic interpolation approach. Thus, for each $a$, either every composition vanishes identically and all nonnegative integers are returns, or one composition has only finitely many zeros and hence permits only finitely many integer returns. The initial times before $a_0$ contribute singletons. Notice that identity of every equation on an analytic arc is enough; no assertion that the arc is algebraically dense is needed.

This is an exact sufficient condition, not an assertion that (2) is always available. Bell--Ghioca--Tucker prove the conjecture for \emph{\'etale} endomorphisms using $p$-adic and algebraic methods [E1]. Their theorem contains a substantial construction step. Simply embedding coefficients into a $p$-adic field and observing that the iterates are defined does not construct a single analytic function of the iterate number. The equations $F(m+1)=f^q(F(m))$ also impose compatibility on an infinite set of arguments, and interpolation by arbitrary functions would not yield the zero theorem.

For a concrete warning, consider $f(z)=z^2$ near a point reducing to $0$ modulo $p$. Its derivative reduces to zero there. Every forward iterate remains in that residue disc, so merely replacing $f$ by an iterate does not make the derivative invertible on this orbit. This example itself has easy return sets in dimension one; it is a countertest to the proposed repair of the method, not a counterexample to dynamical Mordell--Lang. An argument that needs an invertible local derivative cannot silently apply the \emph{\'etale} theorem to this disc.

The inherited 2026 source [S1] concerns split self-maps of an affine curve times a projective curve, with its stated field assumptions. That theorem is valuable additional evidence but does not assert the universal morphism statement printed above. Neither ``split'' nor the special product geometry is automatic for a general quasiprojective variety.

\paragraph{Why the characteristic restriction is structural.}
Over $K=\mathbb F_p(t)$, take
\[
 f(u,v)=(tu,(1-t)v),\qquad x=(1,1),\qquad V:\ u+v=1.
\]
Then $n$ is a return exactly when $t^n+(1-t)^n=1$. Every $n=p^r$ works by Frobenius. Conversely write $n=p^r m>0$ with $p\nmid m$. If $m>1$, the coefficient of $t^{p^r}$ in $(1-t)^n=(1-t^{p^r})^m$ is $-m\ne0$, while neither $t^n$ nor the constant $1$ cancels it. Hence $m=1$. Time zero is not a return. The return set is precisely $\{1,p,p^2,\ldots\}$.

That set is not a finite union of singletons and infinite arithmetic progressions. It is infinite, so such a union would contain an infinite progression; but its successive gaps are unbounded and it contains no infinite progression, since an infinite progression would give a bounded gap between consecutive powers containing its consecutive elements. This explicit failure explains why a characteristic-free polynomial argument cannot prove the requested statement. It also checks that a sparse infinite set need not have the prescribed form: density zero does not imply finiteness.

\paragraph{Verification, first gap, and outcome.}
The accompanying exact checks evaluate the affine example and test the binomial coefficient characterization in small characteristics; these are checks of formulas, not evidence for unrestricted orbit interpolation. The proved results are the quasi-unipotent affine case, the finite-orbit case, and the analytic sufficient condition (2). The first missing step in extending this route is an exact replacement for (2) on arbitrary infinite orbits meeting ramification, or another argument that controls their return times with the same strength. No such replacement is established here. The universal conjecture remains unresolved by this attempt.

\subsection{Sources}
\begin{itemize}
\item[S1]Dynamical Mordell–Lang conjecture for split self-maps of affine curve times projective curve; arXiv2602.08608v2,16April2026.
\url{https://arxiv.org/pdf/2602.08608v2}.
\textit{Formulation source.}
\item[E1]J. P. Bell, D. Ghioca and T. J. Tucker, \textit{The dynamical Mordell--Lang problem for \'etale maps}, American Journal of Mathematics 132 (2010), 1655--1675. The theorem is for \'etale endomorphisms; the analytic criterion in this attempt is stated with its own explicit assumptions.
\url{https://arxiv.org/abs/0808.3266}.

\end{itemize}
\end{document}
